% ============================================================
\chapter{Poisson Process}
\label{ch:poisson}
% ============================================================

How many calls arrive at a call centre between 2pm and 3pm? How many radioactive decays occur in one second? How many customers enter a shop in an hour? In each of these scenarios, events occur randomly in time, independently and at a constant average rate. The process that models exactly this situation is the \emph{Poisson process}, one of the most fundamental objects in probability theory. Discovered by Sim\'eon Denis Poisson in 1837 as a probability distribution, it was elevated to the status of a stochastic process by the work of Palm, Feller, and Khintchine in the 1940s.

\section{Axiomatic Definition}

\begin{definition}[Poisson Process]
\label{def:poisson}
A process $(N_t)_{t \geq 0}$ with values in $\N$ is a \emph{Poisson process}
with intensity $\lambda > 0$ if:
\begin{enumerate}
  \item $N_0 = 0$,
  \item $(N_t)$ has \emph{independent increments}: for all
    $0 \leq t_1 < t_2 < \cdots < t_n$, the r.v.\ $N_{t_2} - N_{t_1},
    N_{t_3} - N_{t_2}, \ldots, N_{t_n} - N_{t_{n-1}}$ are independent,
  \item $(N_t)$ has \emph{stationary increments}:
    $N_{t+s} - N_s \sim \text{Poisson}(\lambda t)$ for all $s, t \geq 0$,
  \item the sample paths of $(N_t)$ are c\`adl\`ag (right-continuous, increasing,
    with jumps of size $1$).
\end{enumerate}
\end{definition}

\begin{remark}
Conditions (1)--(3) suffice to characterise the law of the process.
Condition (4) concerns path regularity and can be obtained by construction.
\end{remark}

\begin{theorem}[Existence]
A Poisson process with intensity $\lambda$ exists for every $\lambda > 0$.
\end{theorem}

\begin{proof}[Proof sketch]
Construct $(N_t)$ from i.i.d.\ exponential inter-arrival times: let
$(T_n)_{n \geq 1}$ be i.i.d.\ $\mathcal{E}(\lambda)$. Set
$S_n = T_1 + \cdots + T_n$ (the $n$-th arrival time) and
$N_t = \max\{n \geq 0 : S_n \leq t\}$. One verifies that $(N_t)$ satisfies
the axioms of Definition~\ref{def:poisson}.
\end{proof}

\section{Fundamental Properties}

\begin{proposition}[Moments]
For all $t \geq 0$:
\[
  \E[N_t] = \lambda t, \quad \Var(N_t) = \lambda t.
\]
\end{proposition}

\begin{proposition}[Inter-Arrival Times]
The times between consecutive jumps $T_n = S_n - S_{n-1}$ (with $S_0 = 0$) are
i.i.d.\ $\mathcal{E}(\lambda)$.
\end{proposition}

\begin{proposition}[Arrival Times]
The $n$-th arrival time $S_n = T_1 + \cdots + T_n$ has a Gamma distribution:
$S_n \sim \text{Gamma}(n, \lambda)$, with density:
\[
  f_{S_n}(t) = \frac{\lambda^n t^{n-1}}{(n-1)!} e^{-\lambda t}, \quad t > 0.
\]
\end{proposition}

\begin{keyformulas}[title=\textbf{Poisson Process Properties}]
\begin{align*}
  N_{t+s} - N_s &\sim \text{Poisson}(\lambda t) \\
  \E[N_t] &= \lambda t, \quad \Var(N_t) = \lambda t \\
  T_n &\sim \mathcal{E}(\lambda) \text{ (i.i.d.)} \\
  S_n &\sim \text{Gamma}(n, \lambda) \\
  \PP(N_t = k) &= \frac{(\lambda t)^k}{k!} e^{-\lambda t}
\end{align*}
\end{keyformulas}

\section{Construction as a CTMC}

\begin{remark}
The Poisson process is a special case of a pure birth process: it is a CTMC on
$\N$ with $q_{n,n+1} = \lambda$ for all $n$ and $q_{ij} = 0$ for $j \neq i + 1$.
The generator is:
\[
  (Qf)(n) = \lambda [f(n+1) - f(n)].
\]
\end{remark}

\section{Markov Property and Memorylessness}

\begin{theorem}[Strong Markov Property]
The Poisson process satisfies the strong Markov property: for every a.s.\ finite
stopping time $\tau$, the process $(N_{\tau+t} - N_{\tau})_{t \geq 0}$ is a
Poisson process with intensity $\lambda$, independent of $\mathcal{F}_{\tau}$.
\end{theorem}

\begin{proposition}[Memorylessness of the Exponential]
The exponential distribution is the unique continuous distribution with the
memoryless property: if $T \sim \mathcal{E}(\lambda)$, then
\[
  \PP(T > t + s \mid T > t) = \PP(T > s) = e^{-\lambda s}, \quad \forall s, t \geq 0.
\]
\end{proposition}

\section{Operations on the Poisson Process}

\subsection{Superposition}

\begin{theorem}[Superposition]
\label{thm:superposition}
Let $(N_t^{(1)})$ and $(N_t^{(2)})$ be independent Poisson processes with
intensities $\lambda_1$ and $\lambda_2$. Then $(N_t^{(1)} + N_t^{(2)})$ is a
Poisson process with intensity $\lambda_1 + \lambda_2$.
\end{theorem}

\subsection{Thinning}

\begin{theorem}[Thinning]
\label{thm:thinning}
Let $(N_t)$ be a Poisson process with intensity $\lambda$. Each jump is
independently classified as type~1 with probability $p$ and type~2 with
probability $1-p$. Then the counting processes $N_t^{(1)}$ and $N_t^{(2)}$
of types~1 and~2 are independent Poisson processes with intensities $\lambda p$
and $\lambda(1-p)$.
\end{theorem}

\begin{proof}
For any $t$, conditionally on $N_t = n$, the $n$ jumps are independently
classified. Hence $N_t^{(1)} \mid (N_t = n) \sim \text{Bin}(n, p)$. We compute:
\[
  \PP(N_t^{(1)} = k) = \sum_{n=k}^{\infty} \binom{n}{k} p^k (1-p)^{n-k}
  \frac{(\lambda t)^n}{n!} e^{-\lambda t} = \frac{(\lambda pt)^k}{k!} e^{-\lambda pt}.
\]
Independence of $N_t^{(1)}$ and $N_t^{(2)}$ follows similarly.
\end{proof}

\subsection{Order Statistics}

\begin{theorem}[Order Statistics]
\label{thm:order_stats}
Conditionally on $N_t = n$, the $n$ jump times $(S_1, \ldots, S_n)$ have the
same distribution as the order statistics of $n$ i.i.d.\ uniform random variables
on $[0, t]$.
\end{theorem}

\section{Compound Poisson Process}

\begin{definition}[Compound Poisson Process]
Let $(N_t)_{t \geq 0}$ be a Poisson process with intensity $\lambda$ and
$(Y_k)_{k \geq 1}$ an i.i.d.\ sequence independent of $(N_t)$. The
\emph{compound Poisson process} is
\[
  Z_t = \sum_{k=1}^{N_t} Y_k, \quad t \geq 0.
\]
\end{definition}

\begin{proposition}[Moments of the Compound Process]
\[
  \E[Z_t] = \lambda t \, \E[Y_1], \quad \Var(Z_t) = \lambda t \, \E[Y_1^2].
\]
\end{proposition}

\begin{example}
An insurance claims model: claims arrive according to a Poisson process with
rate $\lambda$, each claim having a random amount $Y_k$ (i.i.d.). The total
claims amount up to time $t$ is $Z_t = \sum_{k=1}^{N_t} Y_k$.
\end{example}

\section{Non-Homogeneous Poisson Process}

\begin{definition}[Non-Homogeneous Poisson Process]
A process $(N_t)_{t \geq 0}$ is a \emph{non-homogeneous Poisson process} with
\emph{intensity function} $\lambda(t) \geq 0$ if:
\begin{enumerate}
  \item $N_0 = 0$,
  \item $(N_t)$ has independent increments,
  \item $N_{t+s} - N_s \sim \text{Poisson}\bigl(\int_s^{t+s} \lambda(u) \, du\bigr)$.
\end{enumerate}
\end{definition}

\begin{proposition}[Time Change]
Let $\Lambda(t) = \int_0^t \lambda(u) \, du$ be the \emph{cumulative intensity
function}. If $(M_s)_{s \geq 0}$ is a standard Poisson process (intensity $1$),
then $N_t = M_{\Lambda(t)}$ is a non-homogeneous Poisson process with intensity
function $\lambda(t)$.
\end{proposition}

\section{The Poisson Process as a Martingale}

\begin{proposition}[Compensated Martingale]
\label{prop:poisson_martingale}
If $(N_t)$ is a Poisson process with intensity $\lambda$, then:
\begin{enumerate}
  \item $M_t = N_t - \lambda t$ is a martingale (the \emph{compensated Poisson process}).
  \item $M_t^2 - \lambda t$ is a martingale.
  \item $e^{\theta N_t - \lambda t(e^{\theta} - 1)}$ is a martingale for all $\theta \in \R$.
\end{enumerate}
\end{proposition}

\begin{proof}
For item (1):
\begin{align*}
  \E[M_{t+s} - M_s \mid \mathcal{F}_s]
  &= \E[N_{t+s} - N_s - \lambda t \mid \mathcal{F}_s] \\
  &= \E[N_{t+s} - N_s] - \lambda t = \lambda t - \lambda t = 0.
\end{align*}
\end{proof}

\section{Exercises}

\begin{exercise}
Show that if $(N_t)$ is a Poisson process with intensity $\lambda$, then for
fixed $s < t$, $\E[N_s \mid N_t] = N_t \cdot s/t$.
\end{exercise}

\begin{exercise}
Telephone calls arrive according to a Poisson process with rate $\lambda = 5$
per hour. Compute the probability of having at least $3$ calls in $30$ minutes.
\end{exercise}

\begin{exercise}
Let $(N_t)$ be a Poisson process with intensity $\lambda$ and $T > 0$. Show
that $\E[S_1 \mid N_T = n] = T/(n+1)$ where $S_1$ is the first jump time.
\end{exercise}

\begin{exercise}
Let $Z_t = \sum_{k=1}^{N_t} Y_k$ be a compound Poisson process with
$Y_k \sim \mathcal{E}(\mu)$ and $N_t$ Poisson with intensity $\lambda$. Compute
$\E[e^{-\alpha Z_t}]$ for $\alpha > 0$.
\end{exercise}

\begin{exercise}
Prove the superposition theorem (Theorem~\ref{thm:superposition}) using
probability generating functions.
\end{exercise}

\begin{exercise}
Buses arrive at a stop according to a Poisson process with intensity $\lambda$.
A passenger arrives at a deterministic time $t_0 > 0$. Show that the passenger's
waiting time is exponential $\mathcal{E}(\lambda)$, independent of the past.
\end{exercise}
